\chapter{Banach 空间}
\label{chap:banach-spaces}

连续函数的一致范数、可积函数的 \(L^p\) 范数和向量的欧氏范数，都能度量
加法和数乘下的误差。然而同一个代数空间可以容纳不同的范数，某个范数下
存在的极限，在另一个范数下未必存在。本章把这些例子中已经反复使用的
性质整理为赋范空间的语言，再讨论子空间、商空间和线性算子的完备性。
标量域始终记为 \(\mathbb K\)，取 \(\R\) 或 \(\C\)；线性一律相对于选定标量域。

\section{赋范空间}

向量空间 \(X\) 上的范数是满足正定性、绝对齐次性和三角不等式的映射
\(\|\cdot\|:X\to[0,\infty)\)。它诱导平移不变距离
\(d(x,y)=\|x-y\|\)，且
\[
 |\|x\|-\|y\||\le\|x-y\|,\qquad
 \|\alpha x-\beta y\|\le|\alpha|\|x-y\|+|\alpha-\beta|\|y\|.
\]
后一式说明加法、数乘在这个拓扑下连续。范数不只指定哪些点接近，也使
代数运算与极限相容。

前文的 \(\R^n\)、\(C(K)\)、\(L^p(\mu)\)（\(1\le p\le\infty\)）都是例子。
当 \(1\le p<\infty\) 时，复值 \(L^p\) 使用
\(\|f\|_p=(\int|f|^p)^{1/p}\)；\(p=\infty\) 时仍用本质上确界。
Hölder、Minkowski
证明只用模的三角不等式，原样成立。完备性证明中的绝对可和子列也适用于
复值函数，或分别对实部、虚部应用已有实值结果。

\section{完备性与级数}

\begin{definition}
范数诱导的度量完备时，称 \(X\) 为 Banach 空间。记闭单位球为
\(B_X=\{x:\|x\|\le1\}\)；开单位球另写 \(B_X^\circ\)。
\end{definition}

\begin{proposition}\label{prop:banach-series-test}
赋范空间完备当且仅当每个满足 \(\sum_n\|x_n\|<\infty\) 的级数
\(\sum_nx_n\) 都在该空间收敛。
\end{proposition}
\begin{proof}
完备时，部分和之差由范数级数尾和控制，故为 Cauchy 列。反向对任意
Cauchy 列 \(y_n\)，抽取子列使
\(\|y_{n_{k+1}}-y_{n_k}\|\le2^{-k}\)。差级数绝对可和，所以
\(y_{n_k}=y_{n_1}+\sum_{j<k}(y_{n_{j+1}}-y_{n_j})\) 收敛。
原 Cauchy 列随之收敛到相同极限。
\end{proof}

这一判据让“误差分次修正”的构造成为合法方法。以后证明开映射定理时，
将先得到一个近似原像，再把剩余误差按几何级数逐次修正；完备性保证这些
修正仍产生空间中的元素。

\section{子空间与有限积}

Banach 空间的线性子空间在限制范数下完备，当且仅当它闭。这是
度量空间中完备子空间判据的直接应用：子空间内的 Cauchy 列在母空间收敛，
闭性保证极限留下；反向由完备子空间在度量空间中闭得到。

有限个 Banach 空间的积取最大范数或求和范数都是 Banach 空间，因为
Cauchy 性和收敛性可逐坐标判断。无限积则需另加有界性条件。例如
\(\ell^\infty\) 是所有有界数列在上确界范数下的空间；所有数列的代数积中
有无界元素，不能不作限制就使用同一个范数。

\section{商空间}

设 \(M\) 为 \(X\) 的线性子空间。代数商 \(X/M\) 的元素为陪集
\([x]=x+M\)，加法、数乘因等价关系相容而良定义。定义
\[
 \|[x]\|_{X/M}=\inf_{m\in M}\|x-m\|=d(x,M).
\]
换代表元只会平移取下确界的集合，故数值不变。给定 \(\varepsilon>0\)，
可取 \(m,n\in M\) 使 \(\|x-m\|<\|[x]\|+\varepsilon\)、
\(\|y-n\|<\|[y]\|+\varepsilon\)。于是
\[
 \|[x+y]\|\le\|(x-m)+(y-n)\|
 <\|[x]\|+\|[y]\|+2\varepsilon.
\]
令 \(\varepsilon\downarrow0\) 得三角不等式。对非零标量 \(\lambda\)，
由 \(M=\lambda M\) 得
\(\inf_{m\in M}\|\lambda x-m\|=|\lambda|\inf_{n\in M}\|x-n\|\)；
零标量的齐次性直接成立。而 \(\|[x]\|=0\) 等价于 \(x\in\overline M\)。
所以它成为范数恰好要求 \(M\) 闭。
若不要求正定性，只保留非负性、绝对齐次性和三角不等式，就称为半范数。
因此 \(M\) 不闭时，上述公式仍给出半范数，但会把某些非零陪集赋值为零。

\begin{theorem}\label{thm:banach-quotient}
若 \(X\) 为 Banach 空间且 \(M\) 为闭线性子空间，则 \(X/M\) 完备。
自然映射 \(q:X\to X/M\) 连续、满射且开。
\end{theorem}
\begin{proof}
自然映射线性、满射，且 \(\|q(x)\|\le\|x\|\)，因而连续。
若商空间级数满足 \(\sum\|u_n\|<\infty\)，可选代表 \(x_n\) 使
\(\|x_n\|\le\|u_n\|+2^{-n}\)。于是 \(\sum x_n\) 在 \(X\) 收敛，
连续映射 \(q\) 将其部分和映到 \(\sum u_n\) 的部分和，故后者收敛。
上一节级数判据给完备性。这里不要求距离的下确界取得。

对 \(r>0\)，若 \(\|u\|<r\)，下确界定义保证有
代表 \(x\) 满足 \(\|x\|<r\)，故 \(q(B_X^\circ(0,r))=B_{X/M}^\circ(0,r)\)。
平移后即说明 \(q\) 为开映射。
\end{proof}

\section{有界线性算子}

线性映射 \(T:X\to Y\) 称为有界，若存在 \(0\le C<\infty\) 使
\(\|Tx\|\le C\|x\|\) 对所有 \(x\) 成立。这里“有界”指把有界集送到有界集，
不表示 \(T(X)\) 整体有界；非零线性映射的像总是无界的。

\section{算子范数}

有界线性映射全体记为 \(\mathcal L(X,Y)\)，定义
\[
 \|T\|=\sup_{\|x\|\le1}\|Tx\|.
\]
若 \(X\ne\{0\}\)，缩放非零向量还得到
\(\|T\|=\sup_{x\ne0}\|Tx\|/\|x\|\)。若 \(X=\{0\}\)，只有零算子，
其范数由单位闭球公式直接给出为零，不使用空集上的比值上确界。
这恰是所有允许常数 \(C\) 的最小值，且使 \(\mathcal L(X,Y)\) 成为赋范空间。
若复合有意义，则 \(\|ST\|\le\|S\|\|T\|\)。
以下涉及 \(C[a,b]\) 时取有限实数 \(a<b\)。对 \(c\in[a,b]\)，
点值映射 \(f\mapsto f(c)\) 范数为一；积分映射
\(f\mapsto\int_a^bf\) 范数为 \(b-a\)。上界由直接估计，下界用常函数 \(1\)。

\section{连续与有界}

\begin{theorem}\label{thm:linear-continuous-bounded}
线性映射 \(T:X\to Y\) 有界、在零点连续、处处连续三者等价。
\end{theorem}
\begin{proof}
有界给出 Lipschitz 估计 \(\|Tx-Ty\|\le C\|x-y\|\)。处处连续当然蕴含
零点连续。反向，取 \(\delta>0\) 使 \(\|z\|<\delta\) 时 \(\|Tz\|<1\)。
对 \(x\ne0\) 应用于 \(z=\delta x/(2\|x\|)\)，得
\(\|Tx\|<2\|x\|/\delta\)。
\end{proof}

微分算子展示了空间选择的重要性：\(Df=f'\) 从 \(C^1[a,b]\) 配上
\(\|f\|_\infty+\|f'\|_\infty\) 到 \(C[a,b]\) 有界；若定义域只用
\(\|f\|_\infty\)，取 \(f_n(x)=\sin(n(x-a))/n\)，其中 \(n\ge1\)。
则 \(\|f_n\|_\infty\le1/n\to0\)，而 \(f_n'(a)=1\)，故
\(\|Df_n\|_\infty\ge1\)。因此 \(D\) 在这个范数下不连续，也不有界。
一个公式是否定义连续算子，不能脱离定义域和值域的范数判断。

\section{算子空间的完备性}

\begin{theorem}\label{thm:operator-space-complete}
若 \(Y\) 为 Banach 空间，\(X\) 只需赋范，则 \(\mathcal L(X,Y)\) 为 Banach 空间。
\end{theorem}
\begin{proof}
若 \(T_n\) 为算子范数 Cauchy 列，对每个 \(x\)，
\(\|T_nx-T_mx\|\le\|T_n-T_m\|\|x\|\)，故 \(T_nx\) 在 \(Y\) 中收敛，
定义 \(Tx=\lim_nT_nx\)。取极限保留线性；因 \(T_n\) 范数有共同上界 \(C\)，
有 \(\|Tx\|\le C\|x\|\)。最后固定 \(n\) 足够大，先在
\(\|(T_n-T_m)x\|\le\varepsilon\|x\|\) 中令 \(m\to\infty\)，
再对单位球取上确界，得 \(\|T_n-T\|\le\varepsilon\)。
\end{proof}

特别地，连续对偶 \(X^*=\mathcal L(X,\mathbb K)\) 总完备，即使 \(X\) 自身
不完备。逐点极限只能给出一个候选算子，最后在单位球上一致的误差估计
才保证算子范数收敛。

\section{维数的作用}

有限维空间的任意范数与坐标范数等价：上界由基向量范数求和，下界在坐标
单位球面上利用连续范数取得正最小值得到。因此有限维子空间闭且完备，
定义在有限维空间上的任何线性映射都连续。无限维情形下这些断言须分别判断。
\chapterref{chap:continuous-function-space}已由 Riesz 引理证明：闭单位球
范数紧致当且仅当空间有限维。本编将用较弱的拓扑寻找另一种紧致性，
但不会把它误写成范数紧致。

\section*{习题}
\addcontentsline{toc}{section}{习题}
\subsection*{A组}
\begin{enumerate}[label=\textbf{\thechapter.\arabic*.}]
\exerciseitem{banach-czero-closed}证明 \(c_0=\{x=(x_n):x_n\to0\}\) 是 \(\ell^\infty\) 的闭子空间。
\exerciseitem{banach-cfinite-incomplete}证明有限支撑数列空间 \(c_{00}\) 在上确界范数下不完备，其完备化为 \(c_0\)。
\exerciseitem{operator-shift-norm}求 \(\ell^p\) 上右移 \(S(x_0,x_1,\ldots)=(0,x_0,x_1,\ldots)\) 的范数，\(1\le p\le\infty\)。
\exerciseitem{operator-diagonal-norm}设 \(1\le p\le\infty\)。证明公式 \(Tx=(a_nx_n)\) 定义从整个 \(\ell^p\) 到自身的有界线性算子，当且仅当 \((a_n)\) 有界，并求其范数。系数无界时，构造 \(x\in\ell^p\) 使 \((a_nx_n)\notin\ell^p\)。
\end{enumerate}
\subsection*{B组}
\begin{enumerate}[label=\textbf{\thechapter.\arabic*.},resume]
\exerciseitem{quotient-nonclosed}若 \(M\) 不闭，证明商公式只是半范数，并举出非零陪集却半范数为零的例子。
\exerciseitem{operator-uniform-not-pointwise}在 \(\ell^2\) 上令 \(P_nx=x_ne_n\)，证明 \(P_nx\to0\) 对每个 \(x\) 成立，但 \(\|P_n\|=1\)。
\exerciseitem{banach-c-one-complete}证明 \(C^1[a,b]\) 在 \(\|f\|_\infty+\|f'\|_\infty\) 下完备。
\exerciseitem{operator-extension-dense}设 \(D\) 是赋范空间 \(X\) 的稠密线性子空间，\(Y\) Banach。证明每个有界线性算子 \(T:D\to Y\) 都有唯一的有界线性延拓到 \(X\)，且延拓与原算子的范数相等。
\exerciseitem{quotient-distance-constant}在 \(C[a,b]\) 中对常函数子空间取商，证明实函数陪集的商范数为
\[
 \frac{\max f-\min f}{2}.
\]
\end{enumerate}
\subsection*{C组}
\begin{enumerate}[label=\textbf{\thechapter.\arabic*.},resume]
\projectexerciseitem{operator-neumann-series}若 \(X\) Banach、\(A\in\mathcal L(X,X)\) 且 \(\|A\|<1\)，证明
\((I-A)^{-1}=\sum_{n=0}^\infty A^n\)，并给出范数估计。
\end{enumerate}
